\documentclass[aps,prl,reprint,amsmath,amssymb]{revtex4-2}

\usepackage{graphicx}
\usepackage{xcolor}
\usepackage{booktabs}
\usepackage{siunitx}
\usepackage{dcolumn}
\usepackage{bm}
\usepackage{braket}
\usepackage{hyperref}

\newcommand{\gvq}[2]{#2}

\begin{document}

\title{Angular Momentum as Information Geometry:\\Isotropic Packing and the Recurring Angular Cross-Section}

\author{J.~Loutey}
\affiliation{Independent Researcher, Kent, Washington}

\date{February 22, 2026}

\begin{abstract}
\textbf{[SYMBOLIC PROOF]} We show that the angular momentum quantum numbers $(\ell, m)$ emerge
from a geometric construction: uniform-density packing of
distinguishable states on concentric isotropic shells in two
dimensions.  Starting from two axioms---discrete binary
distinguishability and geometric indifference (maximum entropy)---and
the construction's structural choices (motivated, not forced), the
construction yields per-shell angular capacities $2k-1$ for shell~$k$, which is
isomorphic to the spherical harmonic basis $Y_{\ell m}$ with
$\ell = k - 1$ and $m = -\ell, \ldots, +\ell$.  \textbf{[OBSERVATION]}
The principal quantum number~$n$ corresponds to a cumulative grouping
of the first~$n$ shells, producing the $n^2$ orbital count, but this
grouping is not forced by the axioms alone.  Similarly, the per-shell
count carries a factor of~2 from the binary initialization
($N_{\mathrm{init}} = 2$), which we read as corresponding to spin
multiplicity (a counting correspondence, not a derivation of spin).  When the packing is interpreted
as a graph, its Laplacian splits into $\ell$-blocks (the edges change
$n$ or $m$, never $\ell$) and so does not converge to the $S^3$
Laplacian---the spinless node set is the exact $S^3$ label set, on which
$-\Delta_{S^3} = L^2 + A^2$ is the diagonal $SO(4)$ Casimir
$n^2-1$---and \textbf{[INTERNAL THEOREM]} its top
eigenvalue saturates the bipartite bound $\lambda_{\max} \to 2 d_{\max} = 8$
(attained in a mid-$\ell$ block, not in the $s$-wave block, whose own
bound is 4; corrected 2026-09-03), so that under a single matched scale
$K_{\mathrm{vac}} = -1/16$ the extremal eigenvalue reaches the hydrogen
ground-state energy---matched, not derived: the content that could have
failed is the spectral bound, not the target the scale is matched to.
The per-shell spectrum $E_n = -Z^2/(2n^2)$ is not a function of the
graph Laplacian.  On the labels it is the function
$E_n = -Z^2/(2(C+1))$ of the diagonal Casimir $C = n^2-1$, i.e.\ the
bound-state Hamiltonian is $H = -\tfrac{Z^2}{2}(1-\Delta_{S^3})^{-1}$,
and that function is fixed by Fock's continuum projection:\ the eigenvalue $2\pi^2/n$ of
the $S^3$ Coulomb kernel forces the focal length $p_0 = Z/n$, and the
energy shell $p_0^2 = -2E$ then gives $E_n$ (Paper~7).  \textbf{[OBSERVATION]} This angular cross-section recurs across the
natural geometries of the Geometric Vacuum series as a basis
label:\ the full $(\ell, m)$ labeling, with $2\ell+1$ states per~$\ell$,
labels the basis on $S^3$ for hydrogen, the per-electron partial waves
of the hyperspherical treatment of helium, the angular channels of the
molecule-frame hyperspherical treatment, and each fiber of the composed
bundles ($\ell$ is a \emph{conserved} quantum number only for the
one-electron atom;\ elsewhere it is a channel label coupled by Gaunt
coefficients);\ in the two-centre prolate-spheroidal geometry only $m$
survives as a label, because the $\eta$-separation replaces $\ell$ by a
separation index that reduces to it only in the united-atom limit.  This
is Paper~0 in the series.
\end{abstract}

\maketitle

%% ================================================================
%%  I.  INTRODUCTION
%% ================================================================
\section{Introduction}
\label{sec:introduction}

The hydrogen atom possesses $n^2$ orbital states per principal quantum
number~$n$, each labeled by angular momentum
$\ell = 0, 1, \ldots, n-1$ and magnetic quantum number
$m = -\ell, \ldots, +\ell$, with a spin multiplicity of~2 giving
$2n^2$ states total.  The standard derivation obtains these labels
from the eigenvalue structure of the angular momentum operators
$\hat{L}^2$ and $\hat{L}_z$ acting on the Hilbert space of the
Schr\"odinger equation.  Here we ask: to what extent can this quantum
number structure be recovered from geometry alone, without invoking
operators, Hilbert spaces, or the eigenvalue equation?

We present a construction based on uniform-density packing of
distinguishable states on concentric isotropic shells in two
dimensions.  The construction rests on two axioms:
(i)~a minimum of two distinguishable states is required to define a
metric (binary distinguishability), and (ii)~in the absence of
external constraints, the distribution of states must maximize
entropy, enforcing isotropy (geometric indifference).  These axioms,
together with the natural choices of concentric shells and
area-proportional state counting, produce the following results at
three distinct levels of strength.

\medskip

\textbf{[SYMBOLIC PROOF]} \noindent\textbf{Strong claim: angular quantum numbers.}  Shell~$k$
has capacity $2k - 1$, which is isomorphic to the spherical harmonic
basis $Y_{\ell m}$ with $\ell = k - 1$.  The angular momentum labels
$(\ell, m)$ and the per-shell degeneracy $2\ell + 1$ emerge directly
from the packing.  This is the central, mathematically rigorous result
of the paper.

\medskip

\textbf{[OBSERVATION]} \noindent\textbf{Structural correspondences.}  The principal quantum
number~$n$ arises when one groups the first $n$ shells cumulatively,
yielding the count $\sum_{k=1}^{n}(2k-1) = n^2$.  This grouping is
natural but not uniquely forced by the axioms.  The factor of~2 in the
per-shell count comes from the binary initialization
($N_{\mathrm{init}} = 2$);\ the compactified packing manifold $S^2$ is
also orientable (two orientations, a $\mathbb{Z}_2$), which offers a
second, heuristic reading of the same factor.  \textbf{[OBSERVATION]}
The identification of this factor of~2 with spin-$\tfrac{1}{2}$
multiplicity is a counting correspondence, not a derivation from the
axioms.

\medskip

\textbf{[OBSERVATION]} \noindent\textbf{Series significance.}  The $(\ell, m)$ angular
sector constructed here recurs inside the natural geometries of the
Geometric Vacuum (GeoVac) series as a basis label.  The full $(\ell, m)$
labeling, with $2\ell+1$ states per~$\ell$, appears on $S^3$ for
hydrogen~\cite{paper7} (where $\ell$ is also conserved), in the
per-electron angular factors of the hyperspherical coordinates for
He~\cite{paper13}, in the angular channels of the molecule-frame
hyperspherical treatment of H$_2$~\cite{paper15}, and in each fiber of
the composed bundles for LiH~\cite{paper17};\ in the multi-electron and
two-centre cases $\ell$ is a channel label coupled by Gaunt coefficients,
not a conserved quantum number.  In the prolate spheroidal coordinates
for H$_2^+$ and H$_2$~\cite{paper11,paper12} only $m$ is a label:\ the
$\eta$-separation replaces $\ell$ by a separation index whose node count
correlates with $\ell$ only in the united-atom limit.  Where it
appears, the full geometry separates into a system-specific radial or
hyperradial part and an angular cross-section whose labeling is the
packing derived here.  This recurrence is the reason the packing
construction is foundational to the series rather than merely a
curiosity about hydrogen.

\medskip

Several elements of this work connect to well-established results.
The counting $\sum_{\ell=0}^{n-1}(2\ell+1) = n^2$ is standard SO(3)
representation theory.  The relationship between hydrogen's
``accidental'' degeneracy and the geometry of~$S^3$ was established by
Fock~\cite{fock1935} and Bargmann~\cite{bargmann1936}.  What is new
here is the specific packing algorithm that produces the $(\ell, m)$
labeling from geometric--combinatorial reasoning independent of the
operator formalism, and the observation that this angular sector
recurs as the angular fiber of the GeoVac natural geometries wherever
their symmetry retains it (fully at the one-centre, molecule-frame
hyperspherical and composed levels; only through $m$ at the two-centre
prolate level).

The paper is organized as follows.
Section~\ref{sec:construction} presents the packing algorithm.
Section~\ref{sec:angular} establishes the isomorphism between shell
capacity and the spherical harmonic basis.
Section~\ref{sec:correspondences} discusses the structural
correspondences for~$n$ and spin.
Section~\ref{sec:cross-section} shows how the angular cross-section
recurs across the GeoVac natural geometries.
Section~\ref{sec:graph} constructs the graph topology and states what
its spectrum does and does not carry.
Sections~\ref{sec:scope} and~\ref{sec:discussion} address scope,
limitations, and context.


%% ================================================================
%%  II.  THE PACKING CONSTRUCTION
%% ================================================================
\section{The Packing Construction}
\label{sec:construction}

\subsection{Axioms}

The construction rests on two axioms, which contain no reference to
quantum mechanics, operators, wavefunctions, or physical observables.

\medskip

\noindent\textbf{Axiom 1 (Binary Distinguishability).}
To define a metric space, a minimum of two distinguishable states
($N_{\mathrm{init}} = 2$) is required to establish a non-zero
fundamental scale~$d_0$.  A single point cannot define a distance; two
points create the first bit of spatial information.

\textit{Note:} The choice $N_{\mathrm{init}} = 2$ is the minimal
non-trivial initialization.  It is the simplest choice consistent with
this axiom, but not the only one; initializing with $N_{\mathrm{init}} = 3$
or more points would define a different fundamental area and modify the
downstream counting.

\medskip

\noindent\textbf{Axiom 2 (Geometric Indifference).}
In the absence of external constraints or privileged directions, the
distribution of states must maximize entropy.  By the Principle of
Indifference, this necessitates isotropy: states populate isotropic
shells (circles) rather than anisotropic tilings (squares, hexagons),
which would imply privileged directions not present in the axioms.

\textit{Note:} Isotropy motivates circular shells over anisotropic
tilings, but uniform-density packing on concentric circles is one of
several ways to implement isotropy---for instance, hexagonal
close-packing achieves higher packing density while preserving
six-fold symmetry.  Concentric shells with area-proportional counting
is the simplest realization consistent with the axioms, not the only
one.

\medskip

These axioms, together with the specific choices of concentric
circular shells, area-proportional counting, and binary initialization,
constitute the full set of assumptions.  The axioms motivate these
choices but do not uniquely determine them.  What is remarkable is that
this particular realization---arguably the simplest one consistent with
the axioms---produces exactly the angular momentum quantum numbers.

\subsection{Algorithm}

\noindent\textbf{Step 1.}  Place $N_{\mathrm{init}} = 2$ points on an
isotropic shell of radius $r_1 = d_0$, separated by the fundamental
distance~$d_0$.

\textit{Rationale:} Two points are the minimum needed to define a
distance (Axiom~1).  Isotropy requires a circular shell (Axiom~2).
This fixes both $d_0$ and the fundamental area per state
$\sigma_0 = \pi d_0^2 / 2$.

\medskip

\noindent\textbf{Step 2.}  Add concentric isotropic shell~$k$ at
radius $r_k = k \cdot d_0$.

\medskip

\noindent\textbf{Step 3.}  Determine the state count~$N_k$ on
shell~$k$ from uniform information density.  The annular area between
shells $k-1$ and $k$ is
\begin{equation}
  A_k = \pi r_k^2 - \pi r_{k-1}^2
      = \pi d_0^2 \bigl[k^2 - (k-1)^2\bigr]
      = \pi d_0^2 (2k - 1) \,.
  \label{eq:annular-area}
\end{equation}
Dividing by the fundamental area gives the state count:
\begin{equation}
  N_k = \frac{A_k}{\sigma_0}
      = \frac{\pi d_0^2 (2k-1)}{\pi d_0^2 / 2}
      = 2(2k - 1) \,.
  \label{eq:Nk}
\end{equation}

\noindent\textbf{Step 4.}  Repeat for $k = 1, 2, 3, \ldots$

\subsection{Resulting structure}

The construction yields the distribution shown in
Table~\ref{tab:packing}.

\begin{table}[h]
\centering
\begin{tabular}{ccc}
\toprule
Shell $k$ & States $N_k = 2(2k-1)$ & Cumulative $\sum_{i=1}^{k} N_i$ \\
\midrule
1 & 2  & 2  \\
2 & 6  & 8  \\
3 & 10 & 18 \\
4 & 14 & 32 \\
5 & 18 & 50 \\
$\vdots$ & $\vdots$ & $\vdots$ \\
$k$ & $2(2k-1)$ & $2k^2$ \\
\bottomrule
\end{tabular}
\caption{State distribution on isotropic shells.  Each shell~$k$
  carries $2(2k-1)$ states; the cumulative count through shell~$k$ is
  $2k^2$.}
\label{tab:packing}
\end{table}

Two features of this table merit immediate comment.  First, the
per-shell count $2(2k-1)$ factorizes as $2 \times (2k-1)$, where the
factor $2k-1$ is an odd integer determined entirely by the area
formula.  Second, the cumulative count $2k^2$ matches the hydrogen
degeneracy $2n^2$ if one identifies $k = n$.  These observations are
analyzed in the next two sections at appropriately different levels of
strength.


%% ================================================================
%%  III.  ANGULAR QUANTUM NUMBERS FROM UNIFORM DENSITY
%% ================================================================
\section{Angular Quantum Numbers from Uniform Density}
\label{sec:angular}

This section contains the central result of the paper: the per-shell
state count from the packing construction is isomorphic to the angular
momentum labeling $(\ell, m)$.

\subsection{The angular capacity of shell $k$}

The factor of $2(2k-1)$ in Eq.~\eqref{eq:Nk} factorizes into a
global factor of~2---whose origin in the initialization condition (and
a heuristic orientability reading) is analyzed in
Sec.~\ref{sec:correspondences}---and an angular capacity per shell:
\begin{equation}
  N_k^{\mathrm{angular}} = 2k - 1
  \label{eq:angular-capacity}
\end{equation}
distinct angular slots.  \textbf{[SYMBOLIC PROOF]} The annular-area derivation of this $2k-1$
angular capacity (equivalently $2\ell+1$ with $\ell = k-1$) is verified
by \texttt{tests/test\_trunk\_qa\_annular.py}.

\subsection{Isomorphism with spherical harmonics}

The spherical harmonics $Y_{\ell m}(\theta, \phi)$ for a given~$\ell$
have $2\ell + 1$ independent components labeled by
$m = -\ell, -\ell+1, \ldots, +\ell$.  Under the identification
\begin{equation}
  \ell = k - 1 \,,
  \label{eq:ell-k}
\end{equation}
the angular capacity of shell~$k$ becomes
\begin{equation}
  N_k^{\mathrm{angular}} = 2(k-1) + 1 = 2\ell + 1 \,,
  \label{eq:2ell-plus-1}
\end{equation}
which is precisely the dimension of the spin-$\ell$ irreducible
representation of SO(3).  The $2k-1$ angular positions on shell~$k$
are in one-to-one correspondence with the magnetic quantum numbers
$m = -\ell, \ldots, +\ell$.

\textbf{[SYMBOLIC PROOF]} The equality of counts need not be checked case by case: it is an
algebraic identity.  The area formula
(Eq.~\ref{eq:annular-area}) forces the annular area of shell~$k$ to
be proportional to $2k-1$, which is the dimension of the
$(k{-}1)$-th irreducible representation of SO(3).  (The identity
establishes equal counts for every $k$;\ it is not a mechanism linking
planar area to SO(3), and the bijection of positions with $m$ values is
a labelling choice.)  The packing construction therefore generates the
complete angular momentum label set
$\{(\ell, m) : 0 \leq \ell, \; -\ell \leq m \leq \ell\}$ as a direct
consequence of uniform density on isotropic shells.

\subsection{What the construction produces}

To be precise about the scope of this result:

\begin{itemize}
\item The construction produces the \emph{labeling} $(\ell, m)$ and
  the \emph{degeneracy} $2\ell + 1$ per shell.  These are the angular
  quantum numbers.

\item It does \emph{not} produce the functional form of the spherical
  harmonics, the eigenvalue $\ell(\ell+1)$ of $\hat{L}^2$, or the
  commutation relations of angular momentum operators.  Those require
  additional structure: the operators themselves (the $SO(3)$ Casimir
  $L^2$, whose eigenvalue is $\ell(\ell+1)$, and the $so(3)$ commutation
  relations) and, for the functional forms, the continuum harmonics
  reached through Fock's projection (Papers~1 and~7).  A diagonal
  Casimir on the labels supplies eigenvalues, not functional forms or
  commutation relations;\ nor does a continuum limit of the graph
  Laplacian supply them, since it does not converge to $\Delta_{S^3}$
  (Sec.~\ref{sec:graph}).

\item The result is \emph{independent} of any specific physical
  system.  The angular labeling emerges from 2D isotropic geometry,
  not from the hydrogen atom.  Hydrogen merely happens to use this
  same angular sector---as do all rotationally symmetric quantum
  systems.
\end{itemize}


%% ================================================================
%%  IV.  STRUCTURAL CORRESPONDENCES
%% ================================================================
\section{Structural Correspondences: $n$, Spin, and $2n^2$}
\label{sec:correspondences}

The angular quantum numbers $(\ell, m)$ emerge directly from the
packing (Sec.~\ref{sec:angular}).  The principal quantum number~$n$
and the spin factor of~2 involve additional identifications that are
natural and mathematically grounded but not uniquely forced by the two
axioms.  We present them here as structural correspondences.

\subsection{The principal quantum number as cumulative grouping}

If one groups the first $n$ shells cumulatively, the total angular
state count is
\begin{equation}
  \sum_{k=1}^{n} (2k - 1) = n^2 \,,
  \label{eq:n-squared}
\end{equation}
which matches the orbital degeneracy of hydrogen's $n$-th principal
shell: $\sum_{\ell=0}^{n-1}(2\ell+1) = n^2$.  Including the spin
factor gives $2n^2$.

This grouping is natural---the shells are concentrically ordered, so
cumulative counting is the simplest aggregation---but the axioms do
not \emph{require} that shells be grouped this way.  The packing
generates an infinite sequence of shells; the decision to collect the
first~$n$ of them into a single principal shell is an additional
structural choice that corresponds to, but is not derived from, the
axioms.

\subsection{The factor of 2 and orientability}

The per-shell count $N_k = 2(2k-1)$ contains a factor of~2 that
traces back to two sources:

\begin{enumerate}
\item \textbf{Initialization.}  Axiom~1 places
  $N_{\mathrm{init}} = 2$ states on the first shell, fixing the
  fundamental area $\sigma_0 = \pi d_0^2 / 2$.  This factor of~2
  propagates through all subsequent shells via the uniform-density
  condition.

\item \textbf{Orientability (a heuristic reading).}  When the flat
  packing plane is conformally compactified to the Riemann sphere
  $S^2 \cong \mathbb{C} \cup \{\infty\} = \mathbb{CP}^1$, the resulting
  manifold is orientable: it admits two orientations (inward and
  outward normals), a $\mathbb{Z}_2$ that can be read as a doubling of
  the state space.  This $\mathbb{Z}_2$ is the orientation class,
  relating O(3) to SO(3);\ it is \emph{not} the spin double cover
  SU(2)$\to$SO(3), which comes from $\pi_1(\mathrm{SO}(3)) = \mathbb{Z}_2$
  (corrected 2026-09-28:\ an earlier sentence said tracking the
  orientation lifts SO(3) to SU(2)).  The arithmetic factor of~2 is the
  initialization above;\ orientability supplies no additional factor.
\end{enumerate}

\textbf{[OBSERVATION]} The identification of this factor of~2 with the
spin-$\tfrac{1}{2}$ degree of freedom is a counting correspondence:\
the initialization doubling and the spinor representation of SU(2)
both produce a factor of~2.  However, the
packing construction does not derive the spin angular momentum algebra
$[\hat{S}_i, \hat{S}_j] = i\epsilon_{ijk}\hat{S}_k$ or the
eigenvalue $s(s+1)$.  The correspondence is between the \emph{counting}
(a factor of~2) and spin \emph{multiplicity} ($2s+1 = 2$ for
$s = 1/2$), not between the construction and the full spinor
formalism.

\subsection{A remark on the axioms}

The construction follows from the two axioms together with several
additional structural choices: concentric shells with linearly spaced
radii ($r_k = k d_0$), area-proportional state counting, and the
specific fundamental area $\sigma_0 = \pi d_0^2 / 2$ fixed by the
initialization.  These choices are natural and mutually
consistent---they are arguably the simplest possibilities compatible
with the axioms---but they are not uniquely determined by the axioms
alone.  The axioms constrain the construction; they do not uniquely
fix every detail of it.

This distinction matters for the epistemic status of the results.  The
angular quantum numbers $(\ell, m)$ follow from any isotropic,
uniform-density, concentric construction with linearly spaced radii---a
robust class (the linear spacing is needed for the $2k-1$ per-shell
capacity; other spacings change it).  The
specific identification of~$n$ and the factor of~2 depend on the
particular choices above.


%% ================================================================
%%  V.  RECURRING ANGULAR CROSS-SECTION
%% ================================================================
\section{The Packing as Recurring Angular Cross-Section}
\label{sec:cross-section}

The packing construction of Sec.~\ref{sec:construction} produces a
recurring $(\ell, m)$ angular labeling.  In the subsequent papers of
the GeoVac series, this angular sector appears as a structural factor
inside the natural geometries used for quantum chemistry calculations,
as a basis label:\ the full $(\ell, m)$ labeling, with $2\ell+1$ states
per~$\ell$, at Levels~1, 3 and~4 and in each Level-5 fiber ($\ell$
conserved only at Level~1, a Gaunt-coupled channel label elsewhere), but
only $m$ at the two-centre Level~2, where the prolate separation replaces
$\ell$.  Where it appears, the full geometry of each system separates into a
system-specific radial (or hyperradial) coordinate and the angular
cross-section derived here.

Table~\ref{tab:cross-section} summarizes the hierarchy.

\begin{table}[h]
\centering
\begin{tabular}{cllll}
\toprule
Level & System & Full Geometry & Angular Cross-Section & Ref.\ \\
\midrule
1 & H (1-center, 1$e$)
  & $S^3$ (Fock projection)
  & $Y_{\ell m}$ on $S^2 \subset S^3$
  & \cite{paper7} \\
2 & H$_2^+$, H$_2$ (2-center, 1--2$e$)
  & Prolate spheroidal
  & $m$ exact; $\ell$-like index from $\eta$-separation
  & \cite{paper11,paper12} \\
3 & He (1-center, 2$e$)
  & Hyperspherical
  & $Y_{\ell m}$ per electron in fiber
  & \cite{paper13} \\
4 & H$_2$ (2-center, 2$e$)$^{\ast}$
  & Mol-frame hyperspherical
  & $Y_{\ell m}$ angular channels
  & \cite{paper15} \\
5 & LiH (core + valence)
  & Composed fiber bundle
  & $Y_{\ell m}$ per electron group
  & \cite{paper17} \\
\bottomrule
\end{tabular}
\caption{\textbf{[OBSERVATION]} The $(\ell, m)$ angular cross-section
  inside each natural geometry of the GeoVac series.  At Levels~1, 3, 4
  and in each Level-5 fiber the coordinate system separates into a
  system-specific radial part and the full $(\ell, m)$ labeling
  constructed in this paper (a basis label;\ conserved only at
  Level~1);\ at the two-centre Level~2 only $m$ is a label.
  $^{\ast}$H$_2$'s natural geometry is Level~2 (prolate spheroidal,
  Paper~12);\ Level~4 is its structural $N$-electron hyperspherical
  treatment.}
\label{tab:cross-section}
\end{table}

\textbf{Level~1: Hydrogen on $S^3$.}
Fock's 1935 stereographic projection~\cite{fock1935} maps the hydrogen
momentum-space problem onto the three-sphere~$S^3$.  The eigenstates
of the Laplace--Beltrami operator on $S^3$ are hyperspherical
harmonics, which decompose as $Y_{n\ell m} = R_{n\ell}(\chi) \,
Y_{\ell m}(\theta, \phi)$ (the radial factor depends on $\ell$,
$R_{n\ell} \propto \sin^\ell\chi\, C^{\ell+1}_{n-\ell-1}(\cos\chi)$).  The angular factor $Y_{\ell m}$ is
precisely the per-shell structure of the packing, with $\ell = k - 1$.
The packing's node set is exactly the label set of these harmonics
(the Peter--Weyl truncation of $L^2(S^3)$), on which
$-\Delta_{S^3} = L^2 + A^2$ acts as the diagonal $SO(4)$ Casimir
$n^2 - 1$;\ Paper~7~\cite{paper7}
verifies Fock's conformal chain from that continuum operator onward
(under it the hydrogen bound-state Hamiltonian is the function
$H = -\tfrac{Z^2}{2}(1-\Delta_{S^3})^{-1}$ of that operator, its scale fixed by the
eigenvalue $2\pi^2/n$ of Fock's $S^3$ Coulomb kernel;\ that
non-definitional step is tested for $n \le 4$, general $n$ being
Fock's).
The graph's edge Laplacian does \emph{not} converge to it
(Sec.~\ref{sec:graph}).

\textbf{Level~2: H$_2^+$ and H$_2$ on prolate spheroids.}
For diatomics, the natural geometry is the prolate
spheroidal coordinate system $(\xi, \eta, \phi)$, in which the
one-electron Schr\"odinger equation separates;\ it is also H$_2$'s
natural geometry, where a prolate-spheroidal CI with exact algebraic
integrals recovers \textbf{[MEASURED]} $99.97\%$ of the H$_2$
dissociation energy ($0.053$~mHa) with explicit $r_{12}$ correlation,
and $99.81\%$ ($0.32$~mHa) with a re-based orbital CI
(Paper~12~\cite{paper12}).  The azimuthal quantum number~$m$
survives exactly as the quantum number of the $\phi$-equation
($e^{im\phi}$), directly inheriting the packing's angular slot
labeling.  The quantum number~$\ell$ is no longer exactly conserved in
the two-center problem: it is replaced by a separation parameter from
the $\eta$-equation that reduces to~$\ell$ in the united-atom limit.
The angular structure is recognizably descended from the $(\ell, m)$
packing---$m$ is exact---but $\ell$ is no longer a label:\ the prolate
$\eta$-separation replaces it by a separation index whose node count
correlates with $\ell$ only in the united-atom limit (the lifting of the
$(2\ell+1)$-fold degeneracy by the axial potential is common to every
two-centre problem, Level~4 included;\ what is specific to Level~2 is
that the separated coordinates carry no $\ell$ label at all).
Paper~11~\cite{paper11} constructs a discrete lattice in these
coordinates.

\textbf{Level~3: Helium in hyperspherical coordinates.}
For two-electron atoms, the six-dimensional configuration space is
parameterized by a hyperradius~$R$ and five angles.  The angular
eigenstates decompose into single-electron angular momenta
$(\ell_1, m_1)$ and $(\ell_2, m_2)$, each carrying the per-shell
structure of the packing.  Paper~13~\cite{paper13} implements this as
a fiber bundle: the base space is the hyperradius, and the fiber is
the angular manifold.

\textbf{Level~4: H$_2$ in molecule-frame hyperspherical coordinates.}
The structural $N$-electron treatment of two-electron diatomics
combines the molecular axis (internuclear separation~$R$) with the
hyperspherical treatment of the two electrons.  The angular channels
are labeled by per-electron $(\ell, m)$ channel labels, coupled by the
two-centre potential, that determine the multichannel expansion.
\textbf{[MEASURED]} Paper~15~\cite{paper15} recovers 96.0\% of the H$_2$ dissociation
energy using this geometry;\ H$_2$'s natural geometry, and its best
result, is the prolate Level~2 above.

\textbf{Level~5: LiH via composed geometries.}
For core--valence systems, the total geometry is a fiber bundle
$G_{\mathrm{total}} = G_{\mathrm{nuc}} \rtimes G_{\mathrm{core}}(R)
\rtimes G_{\mathrm{val}}(R)$, where each electron group occupies its
own natural geometry.  The angular labeling $(\ell, m)$ appears
independently in each fiber.  Paper~17~\cite{paper17} applies this to
LiH and BeH$^+$.

\medskip

\textbf{[OBSERVATION]} The recurrence of the angular sector across the five levels is not
imposed by hand, and it is not uniform:\ $m$ is an exact label wherever
axial symmetry survives;\ $\ell$ is an exact \emph{basis} label (with
$2\ell+1$ states per $\ell$) wherever the separated coordinates carry an
angular $S^2$ factor---Levels~1, 3, 4 and each Level-5 fiber---but a
\emph{conserved} quantum number only where the full SO(3) survives
(Level~1);\ at Level~2 the prolate separation replaces $\ell$ by a
separation index.  The packing construction provides the combinatorial
substrate for this angular factor.

\subsection{The node structure as combinatorial invariant}
\label{sec:combinatorial_invariant}

The preceding survey reveals a structural feature that merits explicit
comment.  \textbf{[OBSERVATION]} The discrete quantum-number set
$\{(\ell, m) : 0 \leq \ell, \; -\ell \leq m \leq \ell\}$---with its
degeneracy $2\ell + 1$ and selection rules---appears inside several
conformally distinct continuum geometries.  The five geometries catalogued in
Table~\ref{tab:cross-section} are different Riemannian manifolds with
different metrics, curvatures, and coordinate singularity structures.
$S^3$ is compact with constant positive curvature; prolate spheroidal
coordinates are non-compact with coordinate singularities at the foci;
the hyperspherical manifold $S^5$ is compact but five-dimensional (for two electrons); the
composed fiber bundle is a product of manifolds with different
dimensionalities at each level.  Levels~1, 3 and~4, and each fiber of
Level~5, share the full $(\ell, m)$ node structure as a basis labeling,
with its $2\ell+1$ count and angular momentum coupling rules ($\ell$
conserved only at Level~1);\ the two-centre Level~2 shares only the $m$
label, its $\eta$-separation index reducing to $\ell$ only in the
united-atom limit.

This invites a comparison with a standard universality argument
familiar from statistical mechanics and lattice field theory, where many
different lattice discretizations (square, triangular, hexagonal) flow
under renormalization to the same continuum field theory, and the
continuum description is usually read as fundamental.  Here the shared
object runs the other way:\ one discrete label structure (fully, or
through $m$ alone at Level~2) is carried by several continuum
geometries.  This does not make either side prior:\ the label set and
the continuum operators are dual descriptions of the same
quantum-number structure, and the observation is about what they share,
not about which is fundamental.

We state this as a mathematical observation, not a metaphysical claim.
The conformal equivalences between each continuum operator and
standard quantum mechanics are not established uniformly across the
hierarchy.  For the $S^3$ embedding the lattice's labels are the
operator's eigenbasis labels, the hydrogen bound-state Hamiltonian is
the function $H = -\tfrac{Z^2}{2}(1-\Delta_{S^3})^{-1}$ of that operator (Fock), and
\textbf{[SYMBOLIC PROOF]} Paper~7~\cite{paper7} provides 18~symbolic checks of Fock's
continuum projection chain (ten of them limiting-case, definitional or
non-discriminating) (a graph-to-continuum
\emph{operator} convergence is not part of it:\ the edge Laplacian does
not converge to $\Delta_{S^3}$, Sec.~\ref{sec:graph}).  At the other
levels only the angular channel structure carries the packing labels
(at Level~2 the prolate separation constants are not labelled by
$\ell$ at all, only by $m$ and a node count), and the subsequent
papers benchmark each geometry in the hierarchy numerically.
\textbf{[OBSERVATION]} Levels~1--4 are the standard continuum problems
in their separated coordinates; the composed Level~5 adds
a Phillips--Kleinman pseudopotential and is an approximation to, not
an equivalent of, the all-electron problem.
Whether one regards the discrete or continuum description as more
fundamental is a choice of interpretation, not a consequence of the
mathematics.

An analogy to isospectral manifolds is apt.  Kac asked whether one
can hear the shape of a drum~\cite{kac1966};\ Milnor~\cite{milnor1964}
had already exhibited isospectral, non-isometric flat tori, and Gordon,
Webb and Wolpert~\cite{gww1992} answered Kac's question negatively for
planar domains.  \textbf{[OBSERVATION]} The present result differs in a
specific sense:\ the geometries in Table~\ref{tab:cross-section} share
a labeled node structure---quantum number assignments and per-shell
state counts, with the $(\ell, m)$ selection rules shared as basis
labels ($m$ in every row; $\ell$ in every row except the two-center
Level~2, where the prolate separation replaces it).  The spectra themselves differ across rows; what is shared is
combinatorial structure, which is richer than a list of eigenvalues
and different in kind from isospectrality.


%% ================================================================
%%  VI.  GRAPH TOPOLOGY AND WHAT ITS SPECTRUM CARRIES
%% ================================================================
\section{Graph Topology and What Its Spectrum Carries}
\label{sec:graph}

\subsection{Vertex set}

The packing defines a graph $G = (V, E)$ with vertex set
\begin{equation}
  V = \bigl\{(n, \ell, m, s) : 1 \leq n \leq n_{\max},\;
      0 \leq \ell < n,\;
      -\ell \leq m \leq \ell,\;
      s \in \{\uparrow, \downarrow\}\bigr\} \,,
  \label{eq:vertex-set}
\end{equation}
with cardinality
$|V| = \sum_{n=1}^{n_{\max}} 2n^2
     = \tfrac{1}{3}\,n_{\max}(n_{\max}+1)(2n_{\max}+1)$
for a lattice truncated at $n_{\max}$: each shell $n$ contributes its
$2n^2$ spin-doubled capacity, summed over $1 \le n \le n_{\max}$. (The
top shell alone holds $2n_{\max}^2$; the cumulative count is the sum.)
\textbf{[SYMBOLIC PROOF]} This spin-doubled vertex count is verified for
all $n_{\max}$ symbolically (\texttt{tests/test\_paper0\_vertex\_count.py})
and by enumeration (\texttt{tests/test\_dirac\_lattice.py}, which works
in the $(n, \kappa, m_j)$ labelling; the spinless production lattice
holds $\sum n^2$ nodes).

\subsection{Edge set}

Edges encode geometrically adjacent transitions:
\begin{itemize}
\item \textbf{Radial edges} ($\hat{T}_{\pm}$): connect
  $(n, \ell, m, s)$ to $(n \pm 1, \ell, m, s)$ at fixed angular
  quantum numbers $(\ell, m)$.  They connect labels related by the
  ladder operators $\hat{T}_{\pm}$;\ they are not physical transitions.

\item \textbf{Angular edges} ($\hat{L}_{\pm}$): connect
  $(n, \ell, m, s)$ to $(n, \ell, m \pm 1, s)$ at fixed $n$ and
  $\ell$.  They connect labels related by the ladder operators
  $\hat{L}_{\pm}$ within a shell;\ they are not physical transitions.
\end{itemize}

The full graph topology may also include edges connecting states with
different~$\ell$ (e.g., $\ell$-changing transitions driven by
non-spherical perturbations).  The production graph has no
$\ell$-changing edges, so its Laplacian decomposes into independent
$\ell$-blocks ($m$ is mixed within a block by the azimuthal edges ---
corrected 2026-09-03 from ``$(\ell, m)$ blocks'').  This
block-diagonality mirrors the $\ell$-conservation of the spherically
symmetric Coulomb Hamiltonian, but it is a property of the chosen edge
set, not a consequence of the graph encoding the Coulomb problem;\ it
is also what prevents the graph Laplacian from converging to
$\Delta_{S^3}$ (below).

\subsection{Spectral bound}

The graph Laplacian $L = D - A$ defines the legacy graph Hamiltonian
(not the physical hydrogen Hamiltonian)
\begin{equation}
  H = K_{\mathrm{vac}} \, L \,,
  \label{eq:graph-hamiltonian}
\end{equation}
where $K_{\mathrm{vac}} = -1/16$ is the kinetic scale factor.  It
matches only the graph's extremal eigenvalue (a band edge,
$\lambda_{\max} \to 8$) to the hydrogen ground-state energy $-1/2$;
it does not carry the rest of the graph spectrum onto the Rydberg
ladder, whose per-shell energies come from Fock's continuum
projection:\ the $S^3$ Coulomb kernel's eigenvalue $2\pi^2/n$ forces
$p_0 = Z/n$, and the energy shell $p_0^2 = -2E$ then gives $E_n$
(Paper~7;\ the label Casimir $C = n^2-1$ alone does not fix the function,
but once Fock's kernel fixes it, $E_n = -Z^2/(2(C+1))$ on the labels).  \textbf{[OBSERVATION]} This constant's value $1/16$ \emph{coincides} with a conformal-geometry
quantity of the Fock projection (Paper~7~\cite{paper7}): the squared
coupling between adjacent $n$-shells in the Gegenbauer eigenbasis of the
$S^3$ Laplacian is $c^2(n,l) = \tfrac{1}{4}[1 - l(l+1)/(n(n+1))]$ (Paper~7,
corrected 2026-09-03 from a printed $1/16$ prefactor), and the quantity
that coincides with $1/16$ is the inverse Fock Jacobian $1/\Omega^4(0)$
with $\Omega(0) = 2$, equivalently $c^2(n,0)/4$.  This numerical coincidence
between the energy-matching scale and the geometric $1/16$ is an
observation, not a derivation (Paper~7): the matching scale and the
conformal quantity are computed from disjoint inputs, with no known
mechanism connecting them.

\textbf{[INTERNAL THEOREM]} The edges change $n$ (radial, $T_\pm$) or $m$
(azimuthal, $L_\pm$) but never $\ell$, so the Laplacian splits into
$\ell$-blocks; the $\ell = 0$ block is a path graph in $n$ whose top
eigenvalue tends to 4, while the full graph's largest Laplacian
eigenvalue converges to $\lambda_{\max} \to 8 = 2 d_{\max}$, twice the
coordination number, attained in a mid-$\ell$ block ($\ell = 11$ at
$n_{\max} = 30$)
(\texttt{tests/test\_trunk\_qa\_kappa.py},
\texttt{tests/test\_paper7\_graph\_convergence.py}; 2026-09-03
correction of an earlier ``within each $(\ell,m)$ sector'' statement).
\textbf{[INTERNAL THEOREM]} The saturation has a closed-form rate.  Because no edge changes $\ell$, the Laplacian is block diagonal in $\ell$ and each block is the grid graph $P_{n_{\max} - \ell} \times P_{2\ell+1}$ (radial $\times$ azimuthal paths), so the whole spectrum is
\textbf{[PRIOR ART]} the textbook Laplacian spectrum of a Cartesian
product of path graphs~\cite{brouwer_haemers2012} (the block identification and the constant
$C_{\mathrm{sat}}$ below are this paper's)
\begin{equation}
  \operatorname{spec}(L) = \Bigl\{\, 2 - 2\cos\frac{j\pi}{n_{\max}-\ell}
  \;+\; 2 - 2\cos\frac{k\pi}{2\ell+1} \,\Bigr\},
  \label{eq:block_spectrum}
\end{equation}
Each block's top eigenvalue is $4 + 2\cos\frac{\pi}{n_{\max}-\ell} + 2\cos\frac{\pi}{2\ell+1}$, so the deficit is
$\min_\ell\bigl[4\sin^2\frac{\pi}{2(n_{\max}-\ell)} + 4\sin^2\frac{\pi}{2(2\ell+1)}\bigr]$;\ minimising its leading form $\pi^2[(n_{\max}-\ell)^{-2} + (2\ell+1)^{-2}]$ under the block constraint gives the constant in closed form,
\begin{equation}
  2 d_{\max} - \lambda_{\max}
  \;=\; \frac{C_{\mathrm{sat}}}{n_{\max}^{2}} + o\bigl(n_{\max}^{-2}\bigr),
  \qquad
  C_{\mathrm{sat}} \;=\; \frac{\pi^{2}}{4}\bigl(2 + 2^{1/3}\bigr)^{2}\bigl(1 + 2^{-2/3}\bigr)
  \;=\; \gvq{saturation_c}{42.7397}\ldots,
  \label{eq:saturation_rate}
\end{equation}
attained at $\ell/n_{\max} \to 1 - 2/(2 + 2^{1/3}) = 0.386488$ (measured argmax $\ell = 11$ at $n_{\max} = 30$, $27$ at $70$).  \textbf{[MEASURED]} Every sample checked lies \emph{below} $C_{\mathrm{sat}}$ (all $591$
integer cutoffs $10 \le n_{\max} \le 600$) --- $C_{\mathrm{sat},n} = 42.61$ at
$n_{\max} = 320$, $42.70$ at $1000$, $42.73$ at $4000$ --- so every sample
checked understates $C_{\mathrm{sat}}$ (a measured one-sided pattern, not proved for all
$n_{\max}$).  The approach is \emph{not} monotone, however:\ there
are $195$ decreasing steps in that range, e.g.\ $C_{\mathrm{sat},20} = 40.7285 >
C_{\mathrm{sat},21} = 40.6470$ (corrected 2026-09-04;\ ``monotone from below'' was asserted
from a doubling grid that happens to be monotone).  It is the one-sided pattern,
not monotonicity, that makes every checked sample an understatement (the
591-cutoff scan, the $C_{\mathrm{sat},20} > C_{\mathrm{sat},21}$ step and the $n_{\max} = 1000$ value
are re-derived but not pinned by a test), and the value printed here before 2026-09-03, $42.6$, was the $n_{\max} = 320$ sample rather than the limit (\texttt{tests/test\_paper1\_block\_spectrum.py}).  
Finally, with $H = K_{\mathrm{vac}} L$ and $K_{\mathrm{vac}}$ fixed by the match
$K_{\mathrm{vac}} = E_{\mathrm{target}}/\lambda_{\max} = -0.5/8$, the
ground-state eigenvalue of~$H$ is then $-0.5$~Ha by construction.  The
convergence of the spectral bound is the content; the reproduction of
the target is not.  \textbf{[WALL]} Nor is that bound an $S^3$ quantity:\
$8 = 2 d_{\max}$ is the band edge of any large two-dimensional grid, and
$C_{\mathrm{sat}}$ is a Dirichlet band edge.  The load-bearing obstruction is the
kernel dimension:\ because the $\ell$-blocks are disconnected, the
Laplacian on the spinless $(n,\ell,m)$ node set has
$\dim\ker L = n_{\max}$ at every cutoff, whereas $\Delta_{S^3}$ has a
single zero mode, so $L$ does \emph{not} converge to $\Delta_{S^3}$,
for the binary production weights and the $1/(n_i n_j)$ option alike
(\texttt{tests/test\_graph\_s3\_nonconvergence.py}).  (Relatedly, the
graph heat kernel $e^{-tL}$ vanishes identically between blocks;\ turning
that into a statement about the continuum heat kernel would need a
node-to-point identification this paper does not make, so it is
corroboration, not the argument.)  The spin-doubled vertex set of
Eq.~(\ref{eq:vertex-set}) is two copies of this node set;\ the kernel
and label statements here are made on the spinless set.  What is exact is the
spinless node set:\ it is the Gel'fand--Tsetlin label set of the
Peter--Weyl truncation of $L^2(S^3)$, on which
$-\Delta_{S^3} = L^2 + A^2$ acts as the diagonal
$SO(4)$ Casimir $n^2 - 1$ at every cutoff (a node weight, not an edge
Laplacian).  The
\textbf{[SYMBOLIC PROOF]} $(\ell, m)$ angular structure is exact at every finite $n_{\max}$---it
is carried by the node labels, not approximated by mesh refinement.

The conformal equivalence between the hydrogen problem and the
Laplace--Beltrami operator on the unit~$S^3$ (via
Fock's stereographic projection, under which the bound-state Hamiltonian
is the function $H = -\tfrac{Z^2}{2}(1-\Delta_{S^3})^{-1}$ of that operator) is developed in Paper~7~\cite{paper7},
\textbf{[SYMBOLIC PROOF]} where an audit of 18~symbolic checks (ten of them limiting-case, definitional or non-discriminating) validates the conformal geometry
of that continuum chain, and an exact kernel-eigenvalue test carries its
non-definitional energy step, $p_0 = Z/n$ (the graph Laplacian is not a discretization of
it; see above).  Paper~1~\cite{paper1} develops the operator algebra
on these labels (the $SU(1,1)\times SU(2)$ ladder operators, exact on
interior shells, with the Rydberg energies imposed through the energy
shell, $p_0 = Z/n$ fixed by Fock's kernel eigenvalue (Paper~7),
rather than output by the graph); dynamics were attempted via
the graph Laplacian in the archived Paper~6~\cite{paper6} (superseded).


%% ================================================================
%%  VII.  SCOPE AND LIMITATIONS
%% ================================================================
\section{Scope and Limitations}
\label{sec:scope}

\subsection{What the construction provides}

\begin{itemize}
\item[$\checkmark$] The angular quantum number labeling $(\ell, m)$
  and the per-shell degeneracy $2\ell + 1$, as a direct consequence of
  uniform-density isotropic packing.

\item[$\checkmark$] A graph topology $(V, E)$ whose vertices carry
  quantum number labels and whose edges connect labels related by
  ladder operators (not physical transitions).

\item[$\checkmark$] The structural correspondences $n \leftrightarrow$
  cumulative grouping and spin multiplicity $\leftrightarrow$
  the initialization factor of~2 (a counting correspondence), with
  their proper epistemic status identified.
\end{itemize}

\subsection{What the construction does not provide}

\begin{itemize}
\item[$\times$] Energy eigenvalues $E_n = -1/(2n^2)$.  These come
  from Fock's continuum projection (the eigenvalue $2\pi^2/n$ of the
  $S^3$ Coulomb kernel forces $p_0 = Z/n$ on the shell where the
  $SO(4)$ Casimir is $n^2 - 1$, and the energy shell $p_0^2 = -2E$
  converts it to $E_n$;\ Paper~7), not from the graph Laplacian, whose spectrum lies in
  $[0, 8]$;\ the matched kinetic scale $K_{\mathrm{vac}}$ fixes only the
  extremal eigenvalue (its value is matched, and its coincidence with
  the Fock-Jacobian factor $1/16$ is an Observation rather than a
  derivation;\ Papers~1, 7).

\item[$\times$] Wavefunctions $R_{n\ell}(r)$ or
  $Y_{\ell m}(\theta, \phi)$ in their functional form.  The packing
  produces the labels, not the functions.

\item[$\times$] Time evolution and dynamics.  These require the
  Hamiltonian as an operator;\ they were attempted via the graph
  Laplacian in the archived Paper~6~\cite{paper6} (superseded).

\item[$\times$] The fine structure constant or any coupling constants.
  The construction is purely kinematic.
\end{itemize}

\subsection{Choices beyond the axioms}

As noted in Sec.~\ref{sec:correspondences}, the construction involves
structural choices---concentric shells, linearly spaced radii
($r_k = k d_0$), area-proportional counting, and the fundamental area
$\sigma_0 = \pi d_0^2 / 2$---that are natural and mutually consistent
but not uniquely determined by the two axioms.  \textbf{[OBSERVATION]} The robustness of the
$(\ell, m)$ result rests on the observation that any isotropic,
uniform-density, concentric packing with linearly spaced radii
($r_k = k d_0$, so the $k$-th annulus has area proportional to
$k^2 - (k-1)^2 = 2k - 1$) produces the angular capacity $2k - 1$ per
shell;\ the linear spacing is required (proved just below) --- other radial spacings change
the per-shell capacity.  \textbf{[SYMBOLIC PROOF]} Conversely, capacity
$2k - 1$ for every $k$ with $r_0 = 0$ forces $r_k = k d_0$:\ the
condition $r_k^2 - r_{k-1}^2 = (2k-1) d_0^2$ telescopes to
$r_k^2 = \sum_{j \le k} (2j-1)\, d_0^2 = k^2 d_0^2$
(\texttt{tests/test\_trunk\_qa\_annular.py}).


%% ================================================================
%%  VIII.  DISCUSSION
%% ================================================================
\section{Discussion}
\label{sec:discussion}

\subsection{Historical context}

The connection between geometry and quantum mechanics has been
recognized since the earliest days of quantum theory.
Bohr~\cite{bohr1913} imposed quantization as a geometric constraint
(angular momentum $L = n\hbar$ on circular orbits, later reinterpreted
by de~Broglie~\cite{debroglie1925} as an integer number of wavelengths around the orbit).
Fock~\cite{fock1935} discovered that the hydrogen momentum-space
problem maps onto $S^3$ via stereographic projection, revealing the
SO(4) dynamical symmetry that explains the ``accidental''
$n^2$~degeneracy.  Bargmann~\cite{bargmann1936} derived the same
symmetry algebraically.  Malkin and Man'ko~\cite{malkin_manko1965} and
Fronsdal~\cite{fronsdal1967} introduced the conformal group $SO(4,2)$ as
hydrogen's dynamical group, and Barut and Kleinert~\cite{barut1967}
extended the symmetry analysis to transition probabilities via
noncompact dynamical groups.  Biedenharn and Louck~\cite{biedenharn1981}
provided a comprehensive treatment of the angular momentum algebra
underlying the $(\ell, m)$ structure.

The present work operates at a different level: rather than analyzing
the symmetries of the Schr\"odinger equation, it constructs the
angular labeling from geometric--combinatorial reasoning that is
logically independent of the operator formalism.  The result is the
same $(\ell, m)$ basis---as it must be, since both approaches reflect
the same underlying rotational symmetry---but obtained by a different
route.

\subsection{The kinematic--dynamic distinction}

The packing construction produces the angular labelling and
degeneracy counting of the \emph{kinematic} structure of quantum
mechanics.  It does not produce the \emph{dynamics}:
the Hamiltonian, the energy spectrum, or time evolution.  This
separation is a feature of the construction, not a limitation.  It
shows that, given the construction's choices (Sec.~\ref{sec:correspondences}),
the angular quantum-number labels follow from geometry and counting,
independently of the specific force law or potential.

In the GeoVac series, dynamics enter through the Coulomb problem's
continuum structure, not through a limit of the graph Laplacian.
Paper~7~\cite{paper7} verifies Fock's conformal equivalence between the
hydrogen problem and the Laplace--Beltrami operator on~$S^3$ (the
bound-state Hamiltonian being $H = -\tfrac{Z^2}{2}(1-\Delta_{S^3})^{-1}$), under
which the $1/r$ Coulomb potential can equivalently be described as the
coordinate distortion of the stereographic projection (a dual
description, not a claim that either reading is prior);\ the
packing's (spinless) node set is the label set of that operator's
harmonics, on which $-\Delta_{S^3} = L^2 + A^2$ acts as the diagonal
$SO(4)$ Casimir $n^2 - 1$.  The graph
Laplacian equipped with $K_{\mathrm{vac}} = -1/16$ does \emph{not}
converge to that operator:\ it splits into $n_{\max}$ disconnected
$\ell$-blocks (\S VI;\ corrected 2026-09-25 from ``indicated
numerically'', itself a 2026-09-04 correction of ``shows'').

\subsection{Scope of the geometric analogy}

The construction shows that information-theoretic reasoning under
geometric constraints \emph{can} reproduce the angular quantum number
structure.  It does not show that this is the \emph{only} way to
obtain these labels, nor does it establish that quantum mechanics
\emph{must} have a geometric origin.  The angular momentum algebra
$[\hat{L}_i, \hat{L}_j] = i\epsilon_{ijk}\hat{L}_k$ is a consequence
of rotational symmetry regardless of whether one arrives at it through
operator theory or geometric packing.  What the construction
demonstrates is that the \emph{state-counting} aspect of angular
momentum---the degeneracy $2\ell + 1$---has a geometric--combinatorial
derivation, given the construction's structural choices.


%% ================================================================
%%  IX.  CONCLUSIONS
%% ================================================================
\section{Conclusions}
\label{sec:conclusions}

The angular momentum quantum numbers $(\ell, m)$ emerge from
uniform-density packing on isotropic shells: shell~$k$ carries
$2k - 1$ angular slots, isomorphic to the spherical harmonic basis
$Y_{\ell m}$ with $\ell = k - 1$.  This result follows from two
axioms (binary distinguishability and geometric indifference) together
with natural structural choices (concentric shells with linearly
spaced radii $r_k = k\,d_0$, uniform density;\ conversely, capacity
$2k - 1$ for every $k$ forces that spacing).
It is independent of the Schr\"odinger equation, operator postulates,
or any specific physical system.

The principal quantum number~$n$ corresponds to cumulative grouping of
the first~$n$ shells, yielding $n^2$ orbital states.  The factor of~2
in the per-shell count, set by the binary initialization, corresponds
in counting to spin multiplicity.
These are structural correspondences---genuine mathematical
observations that match hydrogen's quantum numbers---but they depend
on identifications beyond what the two axioms alone determine.

The $(\ell, m)$ angular sector constructed here recurs as the
angular cross-section inside the natural geometries of the GeoVac
series as a basis label---in full on $S^3$ for
single-electron atoms~\cite{paper7} (where $\ell$ is also conserved),
as Gaunt-coupled channel labels in the hyperspherical and
molecule-frame hyperspherical treatments and in each fiber of the
composed bundles for core--valence molecules~\cite{paper17};\ only
through $m$ in the two-centre prolate geometry, where the separation
replaces $\ell$.  The computational framework of the series
is built on this label set, on Fock's continuum projection, and on the
natural geometries in which each problem separates---not on the graph
Laplacian and its kinetic scale, whose spectrum is a two-dimensional
grid band structure and does not carry the hydrogen ladder (\S VI).
It has been applied to H, He, H$_2^+$, H$_2$, HeH$^+$, LiH, and
BeH$^+$, benchmarked against exact and experimental values.
\textbf{[MEASURED]} H$_2$ reaches $99.97\%$ of $D_e$ ($0.053$~mHa) in
its natural prolate geometry with explicit $r_{12}$ correlation
(Paper~12~\cite{paper12});\ in that same prolate engine the HeH$^+$
equilibrium bond length reaches $-0.5\%$ at $l_{\max} = 3$ (from
$-4.5\%$ at $l_{\max} = 2$;\ Paper~12);\ the LiH
equilibrium geometry is $+0.2\%$ in an all-electron prolate treatment
with an analytic core (Paper~19~\cite{paper19}) and $5.3\%$ in the
composed Level~5 treatment (Paper~17), where BeH$^+$ is markedly worse
($R_{\rm eq}$ 3.25 vs ${\sim}\,2.48$ bohr with the ab-initio
pseudopotential);\ and in the mol-frame hyperspherical treatment
(Paper~15~\cite{paper15}) HeH$^+$ at fixed bond length recovers 93\% of
$D_e$ adiabatically but only ${\sim}\,5\%$ in the consistent
2D-variational solver.  These inherited figures are owned, and
test-backed or not, by the cited papers.


\section*{Acknowledgments}

This work was developed using an AI-augmented workflow, with
scientific direction and quality control by the author and
implementation assistance from Anthropic's Claude.


\begin{thebibliography}{99}

\bibitem{fock1935}
V.~Fock,
``Zur Theorie des Wasserstoffatoms,''
\textit{Z. Phys.} \textbf{98}, 145 (1935).

\bibitem{bargmann1936}
V.~Bargmann,
``Zur Theorie des Wasserstoffatoms: Bemerkungen zur gleichnamigen
Arbeit von V.~Fock,''
\textit{Z. Phys.} \textbf{99}, 576 (1936).

\bibitem{malkin_manko1965}
I.~A. Malkin and V.~I. Man'ko,
``Symmetry of the hydrogen atom,''
JETP Lett. \textbf{2}, 146 (1965).

\bibitem{fronsdal1967}
C.~Fronsdal,
``Infinite multiplets and the hydrogen atom,''
Phys. Rev. \textbf{156}, 1665 (1967).

\bibitem{debroglie1925}
L.~de~Broglie,
``Recherches sur la th\'eorie des quanta,''
Ann. Phys. (Paris) \textbf{3}, 22 (1925).

\bibitem{brouwer_haemers2012}
A.~E. Brouwer and W.~H. Haemers,
\textit{Spectra of Graphs} (Springer, New York, 2012).

\bibitem{barut1967}
A.~O.~Barut and H.~Kleinert,
``Transition Probabilities of the Hydrogen Atom from Noncompact
Dynamical Groups,''
\textit{Phys. Rev.} \textbf{156}, 1541 (1967).

\bibitem{biedenharn1981}
L.~C.~Biedenharn and J.~D.~Louck,
\textit{Angular Momentum in Quantum Physics: Theory and Application}
(Addison-Wesley, Reading, MA, 1981).

\bibitem{bohr1913}
N.~Bohr,
``On the Constitution of Atoms and Molecules,''
\textit{Phil. Mag.} \textbf{26}, 1 (1913).

\bibitem{paper7}
J.~Loutey,
``The Dimensionless Vacuum: Recovering the Schr\"odinger Equation
on a Scale-Invariant Label Lattice,''
GeoVac Paper~7 (2026).

\bibitem{paper11}
J.~Loutey,
``The Molecular Fock Projection: Diatomic Molecules as Prolate
Spheroidal Lattices,''
GeoVac Paper~11 (2026).

\bibitem{paper12}
J.~Loutey,
``Algebraic Two-Electron Integrals on the Prolate Spheroidal
Lattice,''
GeoVac Paper~12 (2026).

\bibitem{paper13}
J.~Loutey,
``The Hyperspherical Lattice: Two-Electron Atoms as Coupled
Channel Graphs,''
GeoVac Paper~13 (2026).

\bibitem{paper15}
J.~Loutey,
``The Level~4 Natural Geometry: Two-Center Two-Electron Molecules
in Molecule-Frame Hyperspherical Coordinates,''
GeoVac Paper~15 (2026).

\bibitem{paper17}
J.~Loutey,
``Composed Natural Geometries for Core-Valence Diatomics: LiH and
BeH$^+$ from Discrete Graph Topology,''
GeoVac Paper~17 (2026).

\bibitem{paper19}
J.~Loutey,
``Toward PK-Free Molecular Hamiltonians: Two-Center Sturmian Integrals
and Coupled Composition,''
GeoVac Paper~19 (2026).

\bibitem{paper1}
J.~Loutey,
``The Geometric Atom: the Hydrogen State Space as a Label Lattice,''
GeoVac Paper~1 (2026).

\bibitem{kac1966}
M.~Kac,
``Can One Hear the Shape of a Drum?,''
\textit{Amer. Math. Monthly} \textbf{73}, 1 (1966).

\bibitem{milnor1964}
J.~Milnor,
``Eigenvalues of the Laplace Operator on Certain Manifolds,''
\textit{Proc. Natl. Acad. Sci. USA} \textbf{51}, 542 (1964).

\bibitem{gww1992}
C.~Gordon, D.~L.~Webb, and S.~Wolpert,
``One Cannot Hear the Shape of a Drum,''
\textit{Bull. Amer. Math. Soc.} \textbf{27}, 134 (1992).
\bibitem{paper6}
J.~Loutey,
``$O(N)$ Quantum Dynamics and Molecular Spectroscopy via Spectral Graph Theory,''
GeoVac Paper~6 (2026).
[Archived: superseded within the series; retained for provenance.]
\bibitem{paper18}
J.~Loutey,
``Spectral-Geometric Exchange Constants: Projecting Discrete Spectra onto Continuous Manifolds,''
GeoVac Paper~18 (2026).
\end{thebibliography}

\end{document}
